\chapter{Results}
\label{chp:results}

In this chapter the results of the project will be presented. This includes measurements of four transmissions of the image. No compression, 1-bit compression, 2-bit compression and 4-bit compression.

To acquire the results, an image was transferred between the two TelosB motes, one as a sender and another as a receiver. These was implemented as described in chapter \ref{chp:impl}

Three different compressions was developed, and each with a different result, as expected. As shown i figure \ref{fig:one}, \ref{fig:two} and \ref{fig:four}.

\todo{indsæt rette billeder}
\begin{figure}[h]
	\centering
	\begin{subfigure}{0.3\textwidth}
		\includegraphics[width=\textwidth]{oneBit}
		\caption{1-bit compression}
		\label{fig:one}
	\end{subfigure}
	\begin{subfigure}{0.3\textwidth}
		\includegraphics[width=\textwidth]{twoBit}
		\caption{2-bit compression}
		\label{fig:two}
	\end{subfigure}
	\begin{subfigure}{0.3\textwidth}
		\includegraphics[width=\textwidth]{fourbit}
		\caption{4-bit compression}
		\label{fig:four}
	\end{subfigure}
	\caption{Pictures of animals}\label{fig:animals}
\end{figure}


To measure the voltage over the mote and the transmission time, an oscilloscope was used. Figure \ref{fig:receiverresult} shows the measurement of the time for transferring with no compression. 

\myFigure{senderResultconverted}{Measured time for no compression for \emph{Receiver Mote}}{fig:receiverresult}{0.5}


The table below shows all measurements and calculations of the three compressions and the one without compression. The battery voltage for the sender was measured to $3.028V$ and the resistor to $9.93\Omega$. For the receiver the battery voltage was $3.214V$ and the resistor was measured to of $9.98\Omega$.

\vspace{0.5cm}
{\renewcommand{\arraystretch}{1.2}
\begin{tabular}{l l l r r r r}
	Sender           	& $V_{avg}\lbrack mV\rbrack$ & $t\lbrack sec\rbrack$ & $V_{mote}$  & $I_{mote}$ & $W_{mote}\lbrack W\rbrack$ & Energy cons.$\lbrack J \rbrack$\\ \hline
	No compr.	 & 0.194 & 25.2 & 2.834 & 0.020 & 0.055 & 1.395 \\
	1-bit compr. & 0.192 & 24.5 & 2.836 & 0.019 & 0.055 & 1.343 \\
	2-bit compr. & 0.193 & 23   & 2.835 & 0.019 & 0.055 & 1.267 \\
	4-bit compr. & 0.194 & 18.2 & 2.834 & 0.020 & 0.055 & 1.008 \\
	
	&&&&&& \\

	Receiver           	& $V_{avg}\lbrack mV\rbrack$ & $t\lbrack sec\rbrack$ & $V_{mote}$  & $I_{mote}$ & $W_{mote}\lbrack W\rbrack$ & Energy cons.$\lbrack J \rbrack$\\ \hline
	No compr.	 & 0.201 & 25.2 & 3.013 & 0.020 & 0.060 & 1.529 \\
	1-bit compr. & 0.201 & 24.5 & 3.013 & 0.020 & 0.060 & 1.487 \\
	2-bit compr. & 0.2   & 23   & 3.014 & 0.020 & 0.060 & 1.389 \\
	4-bit compr. & 0.2   & 18.2 & 3.014 & 0.020 & 0.060 & 1.099
\end{tabular}

\vspace{0.5cm}

The calculations for $V_{mote}$, $I_{mote}$, $W_{mote}$ and \emph{energy consumption} is listed below.
 
\begin{align*}
	V_{bat}&=V_{bat}-V_{res}\\
	I_{mote}&= \dfrac{V_{avg}}{Res}\\
	W_{mote}&=I_{mote}\cdot V_{mote} \\
	EC&=W_{mote}\cdot t
\end{align*}

The measurement shows a lower energy consumption when applying a higher compression ratio. In the table it is shown how the power used during the transmission, compression and decompression is approximately the same. When applying compression it can though be seen that the transmission time is lower.